%% Plantilla del informe del proceso (LaTeX + Jinja2). La completa app/core/reporte_latex.py.
%% Los delimitadores de Jinja2 están adaptados a LaTeX para no chocar con las llaves: ver
%% reporte_latex._entorno (no se escriben aquí porque Jinja2 los interpretaría).
%% Todos los textos llegan ya escapados para LaTeX desde Python.
%% Fuente: la predeterminada de LaTeX (Latin Modern = Computer Modern vectorial).
\documentclass[11pt,a4paper]{article}
\usepackage[T1]{fontenc}
\usepackage[utf8]{inputenc}
\usepackage{lmodern}
\usepackage{textcomp}
\usepackage[spanish,es-noshorthands,es-tabla]{babel}
\usepackage[a4paper,margin=2.2cm,top=2.6cm,bottom=2.4cm,headheight=14pt]{geometry}
\usepackage{graphicx}
\usepackage{booktabs}
\usepackage{tabularx}
\usepackage{longtable}
\usepackage{array}
\usepackage[table]{xcolor}
\usepackage{amsmath,amssymb}
\usepackage{microtype}
\usepackage{parskip}
\usepackage{fancyhdr}
\usepackage{lastpage}
\usepackage{enumitem}
\usepackage{titlesec}
\usepackage{needspace}
\usepackage[hidelinks]{hyperref}

\definecolor{petroleo}{HTML}{0F4C5C}
\definecolor{tinta}{HTML}{0B1C30}
\definecolor{gris}{HTML}{40484B}
\definecolor{grisclaro}{HTML}{EFF4FF}
\definecolor{verde}{HTML}{15803D}
\definecolor{ambar}{HTML}{B45309}
\definecolor{rojo}{HTML}{B91C1C}
\definecolor{neutro}{HTML}{70787C}
\definecolor{fondoverde}{HTML}{E7F6EC}
\definecolor{fondoambar}{HTML}{FEF3C7}
\definecolor{fondorojo}{HTML}{FDE8E8}
\definecolor{fondoneutro}{HTML}{EEF0F1}

%% Cada título reserva espacio para unas líneas: nunca queda solo al pie de la página.
\titleformat{\section}{\needspace{8\baselineskip}\Large\bfseries\color{petroleo}}{}{0pt}{}
\titleformat{\subsection}{\needspace{6\baselineskip}\large\bfseries\color{tinta}}{}{0pt}{}
%% Las figuras flotan (se ubican donde caben) para no dejar huecos.
\renewcommand{\topfraction}{0.85}
\renewcommand{\textfraction}{0.1}
\renewcommand{\floatpagefraction}{0.75}
\titlespacing*{\section}{0pt}{16pt}{6pt}
\setlist{itemsep=2pt,topsep=2pt}
\renewcommand{\arraystretch}{1.25}

\pagestyle{fancy}
\fancyhf{}
\fancyhead[L]{\small\color{gris}<< proceso >> \textperiodcentered{} << variable >>}
\fancyhead[R]{\small\color{gris}<< fecha >>}
\fancyfoot[L]{\small\color{gris}<< folio >>}
\fancyfoot[R]{\small\color{gris}Página \thepage{} de \pageref{LastPage}}
\renewcommand{\headrulewidth}{0.4pt}

%% \textcolor (no \color) para que el texto arranque en la primera línea de la celda.
\newcommand{\nota}[1]{\textcolor{gris}{\small #1}}

\begin{document}

%% ---------------------------------------------------------------- Portada / encabezado
{\color{petroleo}\rule{\linewidth}{2pt}}\par\vspace{4pt}
{\fontsize{24}{28}\selectfont\bfseries\color{tinta} << titulo >>\par}
\vspace{2pt}
{\large\color{gris} << proceso >> \textperiodcentered{} << variable >>\par}
\vspace{8pt}
\begin{tabularx}{\linewidth}{@{}l X l X@{}}
\nota{Fecha del informe} & << fecha >> & \nota{Folio} & << folio >> \\
\nota{Período analizado} & << periodo >> & \nota{Generado por} & << usuario >> \\
\end{tabularx}

%% ---------------------------------------------------------------- Estado general
\vspace{6pt}
\noindent\colorbox{<< estado.fondo >>}{\begin{minipage}{\dimexpr\linewidth-2\fboxsep\relax}
\vspace{6pt}
{\Large\bfseries\color{<< estado.color >>} << estado.titulo >>\par}
\vspace{2pt}
<< estado.texto >>\par
\vspace{6pt}
\end{minipage}}

<% if hallazgos %>
\section*{En pocas palabras}
\begin{itemize}
<% for h in hallazgos %>  \item << h >>
<% endfor %>\end{itemize}
<% endif %>

<% if indicadores %>
\section*{Indicadores principales}
\begin{longtable}{@{}>{\raggedright\arraybackslash}p{4.0cm} >{\raggedright\arraybackslash}p{4.4cm} >{\raggedright\arraybackslash}p{7.8cm}@{}}
\toprule
\textbf{Indicador} & \textbf{Valor} & \textbf{Qué significa} \\
\midrule
\endhead
<% for fila in indicadores %><< fila[0] >> & \textbf{<< fila[1] >>} & \nota{<< fila[2] >>} \\
<% endfor %>\bottomrule
\end{longtable}
<% endif %>

%% ---------------------------------------------------------------- Comportamiento del proceso
<% if carta %>
\section*{Cómo se comportó el proceso}
<< carta.intro >>

\begin{figure}[htbp]
\centering
\includegraphics[width=\linewidth]{carta.png}\par
\begin{minipage}{\linewidth}
\nota{Cómo leer la gráfica: cada punto es << carta.punto >>. La franja sombreada, entre las líneas rojas,
es la zona normal del proceso. <% if carta.con_espec %>Las líneas punteadas grises son los valores que
acepta el cliente. <% endif %>Los puntos rojos marcados con «!» son muestras fuera de lo normal.}
\end{minipage}
\end{figure}

<% if carta.revisar %>
\subsection*{Muestras fuera de lo normal}
\begin{longtable}{@{}l l l >{\raggedright\arraybackslash}p{8.2cm}@{}}
\toprule
\textbf{Muestra} & \textbf{Fecha y hora} & \textbf{Valor} & \textbf{Qué se observó} \\
\midrule
\endhead
<% for r in carta.revisar %><< r[0] >> & << r[1] >> & << r[2] >> & << r[3] >> \\
<% endfor %>\bottomrule
\end{longtable}
<% else %>
\nota{No hubo muestras fuera de lo normal en el período analizado.}
<% endif %>
<% if carta.excluidas %>

\nota{<< carta.excluidas >>}
<% endif %>
<% endif %>

%% ---------------------------------------------------------------- Cumplimiento de la especificación
<% if capacidad %>
\section*{Cumplimiento de lo que acepta el cliente}
<< capacidad.intro >>

\begin{figure}[htbp]
\centering
\includegraphics[width=\linewidth]{histograma.png}
\end{figure}

\begin{tabularx}{\linewidth}{@{}X r r r@{}}
\toprule
\textbf{Dónde} & \textbf{Esperado} & \textbf{Piezas por millón} & \textbf{Observado} \\
\midrule
<% for fila in capacidad.tabla %><< fila[0] >> & << fila[1] >> & << fila[2] >> & << fila[3] >> \\
<% endfor %>\bottomrule
\end{tabularx}

\nota{«Esperado» se calcula con el promedio y la variabilidad del proceso suponiendo una distribución
normal; «Observado» es lo que efectivamente salió en las muestras.}

\subsection*{Índices de capacidad}
\begin{tabularx}{\linewidth}{@{}>{\raggedright\arraybackslash}p{4.6cm} >{\raggedright\arraybackslash}p{2.2cm} >{\raggedright\arraybackslash}X@{}}
\toprule
\textbf{Índice} & \textbf{Valor} & \textbf{Qué significa} \\
\midrule
<% for fila in capacidad.indices %><< fila[0] >> & \textbf{<< fila[1] >>} & \nota{<< fila[2] >>} \\
<% endfor %>\bottomrule
\end{tabularx}
<% endif %>

%% ---------------------------------------------------------------- Supuestos
<% if supuestos %>
\section*{Verificación de los datos}
<< supuestos.intro >>

\begin{tabularx}{\linewidth}{@{}>{\raggedright\arraybackslash}p{5.4cm} >{\raggedright\arraybackslash}p{3.0cm} >{\raggedright\arraybackslash}X@{}}
\toprule
\textbf{Pregunta} & \textbf{Resultado} & \textbf{Detalle técnico} \\
\midrule
<% for fila in supuestos.filas %><< fila[0] >> & \textbf{<< fila[1] >>} & \nota{<< fila[2] >>} \\
<% endfor %>\bottomrule
\end{tabularx}
<% endif %>

%% ---------------------------------------------------------------- Plan de muestreo
<% if plan %>
\section*{Plan de muestreo}
<< plan.intro >>
\begin{itemize}
<% for t in plan.textos %>  \item << t >>
<% endfor %>\end{itemize}
<% endif %>

%% ---------------------------------------------------------------- Alertas
<% if alertas is not none %>
\section*{Alertas registradas}
<% if alertas %>
\begin{longtable}{@{}>{\raggedright\arraybackslash}p{2.3cm} l >{\raggedright\arraybackslash}p{6.0cm} >{\raggedright\arraybackslash}p{4.4cm}@{}}
\toprule
\textbf{Fecha} & \textbf{Muestra} & \textbf{Qué pasó} & \textbf{Estado y anotaciones} \\
\midrule
\endhead
<% for a in alertas %><< a[0] >> & << a[1] >> & << a[2] >> & << a[3] >> \\
<% endfor %>\bottomrule
\end{longtable}
<% else %>
No se registraron alertas.
<% endif %>
<% endif %>

%% ---------------------------------------------------------------- Lotes
<% if lotes %>
\section*{Inspección de lotes (muestreo de aceptación)}
\begin{longtable}{@{}l l r r r r l@{}}
\toprule
\textbf{Lote} & \textbf{Fecha} & \textbf{Tamaño} & \textbf{Revisadas} & \textbf{Acepta hasta} & \textbf{Defectuosas} & \textbf{Dictamen} \\
\midrule
\endhead
<% for l in lotes %><< l[0] >> & << l[1] >> & << l[2] >> & << l[3] >> & << l[4] >> & << l[5] >> & << l[6] >> \\
<% endfor %>\bottomrule
\end{longtable}
<% endif %>

%% ---------------------------------------------------------------- Historial
<% if historial is not none %>
\section*{Historial de cambios}
<% if historial %>
\begin{longtable}{@{}>{\raggedright\arraybackslash}p{2.3cm} >{\raggedright\arraybackslash}p{2.4cm} >{\raggedright\arraybackslash}p{2.8cm} >{\raggedright\arraybackslash}p{3.4cm} >{\raggedright\arraybackslash}p{4.0cm}@{}}
\toprule
\textbf{Fecha} & \textbf{Quién} & \textbf{Qué cambió} & \textbf{Antes $\rightarrow$ después} & \textbf{Motivo} \\
\midrule
\endhead
<% for h in historial %><< h[0] >> & << h[1] >> & << h[2] >> & << h[3] >> & << h[4] >> \\
<% endfor %>\bottomrule
\end{longtable}
<% else %>
No hay cambios registrados.
<% endif %>
<% endif %>

%% ---------------------------------------------------------------- Anexos
<% if tecnico %>
\section*{Anexo: datos técnicos}
\begin{tabularx}{\linewidth}{@{}X r r r@{}}
\toprule
\textbf{Gráfica} & \textbf{Límite inferior} & \textbf{Línea central} & \textbf{Límite superior} \\
\midrule
<% for fila in tecnico.limites %><< fila[0] >> & << fila[1] >> & << fila[2] >> & << fila[3] >> \\
<% endfor %>\bottomrule
\end{tabularx}

\nota{<< tecnico.nota >>}
<% endif %>

\section*{Glosario}
\begin{description}[style=nextline,leftmargin=0pt,font=\normalfont\bfseries]
<% for g in glosario %>  \item[<< g[0] >>] << g[1] >>
<% endfor %>\end{description}

\vspace{10pt}
\nota{Este informe describe los datos registrados en la aplicación; no incluye recomendaciones.
La integridad del archivo puede verificarse en la aplicación (pantalla Reportes) con su huella SHA-256.}

\end{document}
